\chapter{光滑代数簇上微分算子层上的模}

\section{拟凝聚 $\mathcal{D}_X$-模}

设 $X$ 是拓扑空间，$A$ 是 $X$ 上带单位的环层。以 $\mathcal{M}(A)$ 记 $X$ 上 $A$-模层的范畴，它是 Abel 范畴。以 $A=\Gamma(X,A)$ 记 $A$ 的整体截面环，以 $\mathcal{M}(A)$ 记 $A$-模的范畴。则有自然的加性整体截面函子 $\Gamma=\Gamma(X,-):\mathcal{M}(A)\to\mathcal{M}(A)$。此外存在函子 $\operatorname{Hom}_A(A,-)$ 到函子 $\Gamma(X,-)$ 的自然同构，它把态射 $T:A\to V$ 映为 $T(1_X)\in\Gamma(X,V)$。

还可定义局部化函子 $\Delta:\mathcal{M}(A)\to\mathcal{M}(A)$，由
$$\Delta(V)=A\otimes_A V$$
给出。显然 $\Delta$ 是加性函子，且右正合。

由标准论证，对任意 $V\in\mathcal{M}(A)$ 和 $W\in\mathcal{M}(A)$ 有
$$\operatorname{Hom}_A(A\otimes_A V,W)=\operatorname{Hom}_A(V,\operatorname{Hom}_A(A,W)).$$
故
$$\operatorname{Hom}_A(\Delta(V),W)=\operatorname{Hom}_A(V,\Gamma(X,W)),$$
即 $\Delta$ 是整体截面函子 $\Gamma$ 的左伴随函子。

特别地，存在伴随态射 $\varphi$（从恒等函子到 $\Gamma\circ\Delta$）和 $\psi$（从 $\Delta\circ\Gamma$ 到恒等函子）。

现考虑 $X$ 是代数簇、$A=\mathcal{O}_X$ 是 $X$ 上结构层的特殊情形。此时如前以 $R(X)$ 记 $X$ 上正则函数环。

若 $X$ 是仿射的，称 $\mathcal{M}(\mathcal{O}_X)$ 中的 $V$ 是\emph{拟凝聚 $\mathcal{O}_X$-模}，若存在 $R(X)$-模 $V$ 使得 $V\cong\Delta(V)$。

若 $X$ 是任意代数簇，称 $V$ 是拟凝聚 $\mathcal{O}_X$-模，若每点 $x\in X$ 有仿射开邻域 $U$ 使得 $V|_U$ 是拟凝聚 $\mathcal{O}_U$-模。拟凝聚 $\mathcal{O}_X$-模构成 $\mathcal{M}(\mathcal{O}_X)$ 的全子范畴，记为 $\mathcal{M}_{qc}(\mathcal{O}_X)$。可以检验 $\mathcal{M}_{qc}(\mathcal{O}_X)$ 是 Abel 范畴。

显然 $\Delta$ 是从 $\mathcal{M}(R(X))$ 到 $\mathcal{M}_{qc}(\mathcal{O}_X)$ 的函子。若 $X$ 仿射，由 Serre 定理 $\Delta:\mathcal{M}(R(X))\to\mathcal{M}_{qc}(\mathcal{O}_X)$ 是范畴的等价，$\Gamma:\mathcal{M}_{qc}(\mathcal{O}_X)\to\mathcal{M}(R(X))$ 是其拟逆。

设 $\mathcal{D}_X$ 是 $X$ 上微分算子层。则有自然同态 $\iota:\mathcal{O}_X\to\mathcal{D}_X$，它定义从 $\mathcal{M}(\mathcal{D}_X)$ 到 $\mathcal{M}(\mathcal{O}_X)$ 的遗忘函子。称 $\mathcal{D}_X$-模 $V$ 是拟凝聚的，若它是拟凝聚 $\mathcal{O}_X$-模。

以 $\mathcal{M}_{qc}(\mathcal{D}_X)$ 记由拟凝聚 $\mathcal{D}_X$-模构成的 $\mathcal{M}(\mathcal{D}_X)$ 的全子范畴。则 $\mathcal{M}_{qc}(\mathcal{D}_X)$ 是 Abel 范畴。

再设 $X$ 是仿射簇。以 $D_X=\Gamma(X,\mathcal{D}_X)$ 记 $X$ 上整体微分算子环。则由 Serre 定理 $\Gamma:\mathcal{M}_{qc}(\mathcal{D}_X)\to\mathcal{M}(D_X)$ 是正合函子。此外 $\Delta$ 是从 $\mathcal{M}(D_X)$ 到 $\mathcal{M}_{qc}(\mathcal{D}_X)$ 的函子。有如下 Serre 定理的类似物。

\begin{theorem}
设 $X$ 是仿射簇。则 $\Gamma:\mathcal{M}_{qc}(\mathcal{D}_X)\to\mathcal{M}(D_X)$ 是范畴的等价。局部化函子 $\Delta:\mathcal{M}(D_X)\to\mathcal{M}_{qc}(\mathcal{D}_X)$ 是其拟逆。
\end{theorem}

\begin{proof}
设 $V\in\mathcal{M}(D_X)$。则存在 $D_X$-模的正合列 $D_X^{(I)}\to D_X^{(J)}\to V\to 0$，应用 $\Delta$ 后得到 $\mathcal{D}_X$-模的正合列 $\mathcal{D}_X^{(I)}\to\mathcal{D}_X^{(J)}\to\Delta(V)\to 0$。函子 $\Gamma\circ\Delta$ 是 $\mathcal{M}(D_X)$ 到自身的右正合函子。此外对任意 $V\in\mathcal{M}(D_X)$ 有伴随态射 $\varphi_V:V\to\Gamma(X,\Delta(V))$。断言它是同构。显然对任意自由 $D_X$-模 $F$，$\varphi_F:F\to\Gamma(X,\Delta(F))$ 是同构。因此取 $D_X$-模的正合列 $D_X^{(I)}\to D_X^{(J)}\to V\to 0$，得到如下 $D_X$-模的交换图
$$\begin{tikzcd}
D_X^{(I)} \arrow[r] \arrow[d,"\varphi_{D_X^{(I)}}"'] & D_X^{(J)} \arrow[r] \arrow[d,"\varphi_{D_X^{(J)}}"'] & V \arrow[r] \arrow[d,"\varphi_V"] & 0\\
D_X^{(I)} \arrow[r] & D_X^{(J)} \arrow[r] & \Gamma(X,\Delta(V)) \arrow[r] & 0
\end{tikzcd}$$
其行正合且前两个竖直箭头是同构，因此 $\varphi_V$ 是同构。

现考虑另一伴随态射。对任意拟凝聚 $\mathcal{D}_X$-模 $V$，存在自然态射 $\psi_V:\Delta(\Gamma(X,V))\to V$。断言它是同构。考虑拟凝聚 $\mathcal{D}_X$-模的正合列
$$0\longrightarrow K\longrightarrow\Delta(\Gamma(X,V))\xrightarrow{\ \psi_V\ }V\longrightarrow C\longrightarrow 0.$$
由于 $\Gamma(X,-)$ 正合，得到 $D_X$-模的正合列
$$0\to\Gamma(X,K)\to\Gamma(X,\Delta(\Gamma(X,V)))\to\Gamma(X,V)\to\Gamma(X,C)\to 0.$$
由证明的第一部分，中间箭头是同构，故 $\Gamma(X,K)=\Gamma(X,C)=0$。由 Serre 定理最终断定 $K=C=0$。
\end{proof}

回顾 $X$ 上层 $F$ 的\emph{支集} $\operatorname{supp}(F)$ 是使 $F|_U=\{0\}$ 的最大开集 $U$ 的补集，因此层的支集是闭的。设 $U$ 是 $X$ 中开集，$s$ 是 $F$ 在 $U$ 上的局部截面，则 $\operatorname{supp}(s)$ 是使 $s|_V=0$ 的最大开集 $V\subset U$ 的补集。

\begin{lemma}
对任意 $s\in F(U)$，有
$$\operatorname{supp}(s)=\{x\in U\mid s_x\neq 0\}.$$
\end{lemma}

\begin{proof}
显然若 $x\notin\operatorname{supp}(s)$，存在 $x$ 的开邻域 $V$ 使 $s|_V=0$ 且 $s_x=0$。反之若 $s_x=0$，存在 $x$ 的开邻域 $V\subset U$ 使 $s|_V=0$ 且 $x\notin\operatorname{supp}(s)$。
\end{proof}

此外有如下结果。

\begin{proposition}
设 $F$ 是 $X$ 上的层。则 $\operatorname{supp}(F)$ 是集合 $\{x\in X\mid F_x\neq 0\}$ 的闭包。
\end{proposition}

\begin{proof}
显然 $\operatorname{supp}(F)$ 必含其所有截面的支集，故 $\{x\in X\mid F_x\neq 0\}$ 必含于支集中。由于 $F$ 的支集是闭的，它也含 $\{x\in X\mid F_x\neq 0\}$ 的闭包。设 $x$ 在此闭包之外，则存在 $x$ 的开邻域 $U$ 使对所有 $y\in U$ 有 $F_y=0$。但这意味着 $F|_U=0$ 且 $U\cap\operatorname{supp}(F)=\varnothing$，故 $x\notin\operatorname{supp}(F)$。
\end{proof}

设 $X=k^n$，$V$ 是 $D(n)$-模。则 $\mathcal{V}=\Delta(V)$ 是拟凝聚 $\mathcal{D}_X$-模。此外对任意 $x\in X$ 有 $\mathcal{V}_x=V_x$。因此由上述命题，$\operatorname{supp}(\mathcal{V})$ 等于 I.4 意义下 $\operatorname{supp}(V)$ 在 $X$ 中的闭包。

\begin{proposition}
设
$$0\longrightarrow F_1\longrightarrow F_2\longrightarrow F_3\longrightarrow 0$$
是 $X$ 上层短正合列。则
$$\operatorname{supp}(F_2)=\operatorname{supp}(F_1)\cup\operatorname{supp}(F_3).$$
\end{proposition}

\begin{proof}
由假设，对所有 $x\in X$ 序列
$$0\longrightarrow F_{1,x}\longrightarrow F_{2,x}\longrightarrow F_{3,x}\longrightarrow 0$$
正合。因此
$$\{x\in X\mid F_{2,x}\neq 0\}=\{x\in X\mid F_{1,x}\neq 0\}\cup\{x\in X\mid F_{3,x}\neq 0\}.$$
故取闭包并用上述命题即得结论。
\end{proof}
\begin{annote}
$\mathcal D_X=\varinjlim\mathcal D_{X,\le d}$ 是 $\mathcal O_X$-拟凝聚层；它在光滑点处局部同构于 $\mathcal O_X\otimes\mathrm{Sym}\,\mathcal T_X$ 的完成，但整体上不是 $\mathcal O_X$-凝聚层。讲义此后凡说“$\mathcal D_X$-模”都指左模，右模与双边模的换边靠 $\omega_X$ 与 $\mathcal D_X$ 的自反性。[来源：Stacks Project Tag 01BD（Section 17.10, Quasi-coherent modules）、Tag 01BU（Section 17.12, Coherent modules）；原书 §1, p77–86]
\end{annote}



\section{凝聚 $\mathcal{D}_X$-模}

现设 $X$ 是光滑仿射簇，则 $D_X$ 是诺特环。因此 $\mathcal{M}(\mathcal{D}_X)$ 中由有限生成 $\mathcal{D}_X$-模构成的全子范畴 $\mathcal{M}_{fg}(\mathcal{D}_X)$ 是 Abel 范畴。称 $V$ 是\emph{凝聚 $\mathcal{D}_X$-模}，若 $V\cong\Delta(V_0)$，其中 $V_0$ 是某有限生成 $D_X$-模。以 $\mathcal{M}_{coh}(\mathcal{D}_X)$ 记由凝聚 $\mathcal{D}_X$-模构成的 $\mathcal{M}(\mathcal{D}_X)$ 的全子范畴。显然 $\Gamma$ 把 $\mathcal{M}_{coh}(\mathcal{D}_X)$ 映入 $\mathcal{M}_{fg}(\mathcal{D}_X)$，$\Delta$ 把 $\mathcal{M}_{fg}(\mathcal{D}_X)$ 映入 $\mathcal{M}_{coh}(\mathcal{D}_X)$。因此 $\Gamma:\mathcal{M}_{coh}(\mathcal{D}_X)\to\mathcal{M}_{fg}(\mathcal{D}_X)$ 是范畴的等价，$\Delta:\mathcal{M}_{fg}(\mathcal{D}_X)\to\mathcal{M}_{coh}(\mathcal{D}_X)$ 是其拟逆。故此时可把凝聚性看作有限生成性的层化版本。

\begin{lemma}
设 $X$ 是光滑仿射簇，$V$ 是拟凝聚 $\mathcal{D}_X$-模。则下列条件等价：
\begin{enumerate}[label=(\roman*)]
  \item $V$ 是凝聚 $\mathcal{D}_X$-模；
  \item 对任意 $x\in X$ 存在 $x$ 的开邻域 $U$ 和正合列 $\mathcal{D}_U^p\to\mathcal{D}_U^q\to V|_U\to 0$。
\end{enumerate}
\end{lemma}

\begin{proof}
(i)$\Rightarrow$(ii) 设 $V$ 凝聚，则 $V\cong\Delta(V_0)$，其中 $V_0$ 是有限生成 $D_X$-模。由于 $D_X$ 诺特，存在正合列 $D_X^p\to D_X^q\to V_0\to 0$（$p,q\in\mathbb{Z}_+$）。局部化得正合列 $\mathcal{D}_X^p\to\mathcal{D}_X^q\to\Delta(V_0)\to 0$。由于 $V\cong\Delta(V_0)$，存在正合列 $\mathcal{D}_X^p\to\mathcal{D}_X^q\to V\to 0$。故对任意 $x\in X$ 可取 $U=X$。

(ii)$\Rightarrow$(i) 存在 $f\in R(X)$ 使 $f(x)\neq 0$ 且 $X_f\subset U$。故缩小 $U$ 可设为主开集。则 (ii) 蕴含序列 $\mathcal{D}_U^p\to\mathcal{D}_U^q\to\Gamma(U,V)\to 0$ 正合，即 $\Gamma(U,V)$ 是有限生成 $D_U$-模。现 $\Gamma(U,V)=\Gamma(X,V)_f$，且存在 $v_1,\dots,v_n\in\Gamma(X,V)$ 使其在 $U$ 上的限制作为 $D_U$-模生成 $V|_U$。所有这些主开集构成 $X$ 的开覆盖。由于 $X$ 拟紧，可取有限子覆盖，从而存在 $w_1,\dots,w_m\in\Gamma(X,V)$ 使得每个茎 $V_x$ 作为 $\mathcal{D}_{X,x}$-模由它们的像生成。因此有满射 $\mathcal{D}_X^m\to V$，进而有满射 $D_X^m\to\Gamma(X,V)$，$\Gamma(X,V)$ 是有限生成 $D_X$-模。故 $V$ 凝聚。
\end{proof}

设 $X$ 是任意光滑代数簇。称 $X$ 上拟凝聚 $\mathcal{D}_X$-模 $V$ 是凝聚的，若对任意 $x\in X$ 存在 $x$ 的开邻域 $U$ 和正合列 $\mathcal{D}_U^p\to\mathcal{D}_U^q\to V|_U\to 0$。由引理 2.1，此定义与仿射簇情形的先前定义一致。此外前面的结果蕴含如下命题。

\begin{proposition}
设 $V$ 是光滑代数簇 $X$ 上的拟凝聚 $\mathcal{D}_X$-模。则下列条件等价：
\begin{enumerate}[label=(\roman*)]
  \item $V$ 是凝聚 $\mathcal{D}_X$-模；
  \item 对任意仿射开子集 $U\subset X$，限制 $V|_U$ 是凝聚 $\mathcal{D}_U$-模；
  \item 对 $X$ 的仿射开覆盖 $(U_1,\dots,U_n)$，限制 $V|_{U_i}$（$1\le i\le n$）是凝聚 $\mathcal{D}_{U_i}$-模。
\end{enumerate}
\end{proposition}

以 $\mathcal{M}_{coh}(\mathcal{D}_X)$ 记 $\mathcal{M}_{qc}(\mathcal{D}_X)$ 中由凝聚 $\mathcal{D}_X$-模构成的全子范畴。则 $\mathcal{M}_{coh}(\mathcal{D}_X)$ 是 Abel 范畴。

对凝聚 $\mathcal{D}_X$-模可改进前面的支集结果。

\begin{proposition}
设 $V$ 是凝聚 $\mathcal{D}_X$-模。则
$$\operatorname{supp}(V)=\{x\in X\mid V_x\neq 0\}.$$
\end{proposition}

\begin{proof}
由前面的命题，只需证 $\{x\in X\mid V_x\neq 0\}$ 是闭集。设 $y$ 是其闭包中一点，$U$ 是 $y$ 的仿射邻域。则 $V(U)$ 是有限生成 $D_U$-模。设 $s_1,\dots,s_n$ 是 $V(U)$ 中作为 $D_U$-模生成它的截面，则这些截面也作为 $\mathcal{D}_U$-模生成 $V|_U$。令 $Z=\bigcup_{i=1}^n\operatorname{supp}(s_i)$，则 $Z$ 是 $U$ 中含于 $\{x\in X\mid V_x\neq 0\}$ 的闭子集。假设 $y\notin Z$，则存在 $y$ 的开邻域 $W\subset U$ 使 $s_1,\dots,s_n$ 在 $W$ 上消失。由此可知 $V|_W=0$，$y\notin\operatorname{supp}(V)$，与前面的命题矛盾。故 $y\in Z$。
\end{proof}

设 $U$ 是 $X$ 中开集，$F$ 是 $\mathcal{O}_X$-模，$G$ 是 $F|_U$ 的 $\mathcal{O}_U$-子模。以 $\bar G$ 记由
$$\bar G(V)=\{s\in F(V)\mid s|_{V\cap U}\in G(V\cap U)\}$$
定义的 $F$ 的子层。显然 $\bar G$ 是 $\mathcal{O}_X$-子模，称为 $G$ 的\emph{典范扩张}。

\begin{lemma}
设 $F$ 是拟凝聚 $\mathcal{O}_X$-模，$G$ 是 $F|_U$ 的拟凝聚 $\mathcal{O}_U$-子模。则 $G$ 的典范扩张 $\bar G$ 是拟凝聚 $\mathcal{O}_X$-模。
\end{lemma}

\begin{proof}
设 $i$ 是 $U$ 到 $X$ 的自然包含。以 $H$ 记 $F|_U$ 模 $G$ 的商，则 $H$ 是拟凝聚 $\mathcal{O}_U$-模。考虑拟凝聚 $\mathcal{O}_X$-模的自然态射 $\alpha:i_*(F|_U)\to i_*(H)$。它与典范态射 $F\to i_*(F|_U)$ 的复合定义拟凝聚 $\mathcal{O}_X$-模态射 $\varphi:F\to i_*(H)$。故其核拟凝聚，且
\begin{align*}
\ker\varphi(V)&=\{s\in F(V)\mid \varphi_V(s)=0\}\\
&=\{s\in F(V)\mid \alpha_{V\cap U}(s|_{V\cap U})=0\}\\
&=\{s\in F(V)\mid s|_{V\cap U}\in G(V\cap U)\}=\bar G(V).
\end{align*}
\end{proof}

设 $V$ 是拟凝聚 $\mathcal{D}_X$-模，$W$ 是 $V|_U$ 的拟凝聚 $\mathcal{D}_U$-子模。则 $W$ 的典范扩张 $\bar W$ 是 $V$ 的拟凝聚 $\mathcal{D}_X$-子模。

非零拟凝聚 $\mathcal{D}_X$-模 $V$ 称为\emph{不可约的}，若 $V$ 的任意拟凝聚 $\mathcal{D}_X$-子模 $W$ 或为 $\{0\}$ 或等于 $V$。

\begin{lemma}
设 $U$ 是 $X$ 中开集，$V$ 是不可约拟凝聚 $\mathcal{D}_X$-模。则 $V|_U$ 或为不可约拟凝聚 $\mathcal{D}_U$-模，或为零。
\end{lemma}

\begin{proof}
设 $V|_U\neq 0$。设 $W$ 是 $V|_U$ 的拟凝聚 $\mathcal{D}_U$-子模。以 $\bar W$ 记它到 $V$ 的 $\mathcal{D}_X$-子模的典范扩张。由于 $V$ 不可约，$\bar W$ 或为 $V$ 或为 $0$。这意味着 $W$ 或为 $V|_U$ 或为 $0$。
\end{proof}

特别地有如下推论。

\begin{proposition}
设 $V$ 是不可约拟凝聚 $\mathcal{D}_X$-模。则 $V$ 是凝聚的。
\end{proposition}

\begin{proof}
设 $U$ 是 $X$ 中仿射开集。则由前面的引理 $V|_U$ 或不可约或为 $0$。若 $V|_U$ 不可约，由 1.1，$\Gamma(U,V)$ 必为不可约 $D_U$-模。故 $V|_U$ 是凝聚 $\mathcal{D}_U$-模。结论由前面的命题得出。
\end{proof}

\begin{proposition}
设 $V$ 是不可约拟凝聚 $\mathcal{D}_X$-模。则支集 $\operatorname{supp}(V)$ 是 $X$ 的不可约闭子簇。
\end{proposition}

\begin{proof}
由定义 $\operatorname{supp}(V)$ 是 $X$ 的闭子簇。先断言 $\operatorname{supp}(V)$ 连通。假设 $\operatorname{supp}(V)$ 是 $X$ 的两个闭子簇 $Z_1$ 和 $Z_2$ 的不交并，且 $Z_1\neq\varnothing$。令 $U=X-Z_1$，以 $W$ 记 $V|_U$ 的零 $\mathcal{D}_U$-子模的典范扩张。由于 $V$ 不可约，$W$ 或为 $V$ 或为 $0$。取 $x\in Z_1$ 和不与 $Z_2$ 相交的 $x$ 的仿射开邻域 $W'$。则 $V|_{W'}$ 的支集等于 $Z_1\cap W'=(X-U)\cap W'=W'-(W'\cap U)$。另一方面，
$$\Gamma(W',W)=\{s\in\Gamma(W',V)\mid s|_{W'\cap U}=0\}=\Gamma(W',V),$$
故 $W\neq 0$，从而 $W=V$，$V|_U=0$，故 $Z_2=\varnothing$。因此 $\operatorname{supp}(V)$ 连通。

现证 $\operatorname{supp}(V)$ 不可约。反设不然。设 $Z_1$ 是 $\operatorname{supp}(V)$ 的一个不可约分支，$Z_2$ 是其余所有不可约分支的并，则 $\operatorname{supp}(V)=Z_1\cup Z_2$。令 $Z=Z_1\cap Z_2$，$U=X-Z$，则 $V|_U\neq 0$。由前面的引理它是不可约 $\mathcal{D}_U$-模。显然其支集等于 $(Z_1\cup Z_2)-(Z_1\cap Z_2)=(Z_1-(Z_1\cap Z_2))\cup(Z_2-(Z_1\cap Z_2))$。由前述结果此空间必连通，故 $Z_2-(Z_1\cap Z_2)=\varnothing$。由此可知 $Z_2\subset Z_1$，矛盾。
\end{proof}

拟凝聚 $\mathcal{D}_X$-模 $V$ 称为\emph{有限长度的}，若它有由拟凝聚 $\mathcal{D}_X$-模构成的有限递增滤过
$$\{0\}=V_0\subsetneq V_1\subsetneq\cdots\subsetneq V_n=V,$$
使得 $V_p/V_{p-1}$ 是不可约 $\mathcal{D}_X$-模。数 $\ell(V)=n$ 称为 $V$ 的\emph{长度}。显然由长度上的归纳和前面的命题，任意有限长度的拟凝聚 $\mathcal{D}_X$-模是凝聚的。

\begin{lemma}
设 $V$ 是拟凝聚 $\mathcal{D}_X$-模。则下列条件等价：
\begin{enumerate}[label=(\roman*)]
  \item $V$ 具有有限长度；
  \item 对任意开子集 $U\subset X$，限制 $V|_U$ 具有有限长度；
  \item 存在 $X$ 的开覆盖 $(U_i;\ 1\le i\le n)$ 使得 $V|_{U_i}$（$1\le i\le n$）具有有限长度。
\end{enumerate}
\end{lemma}

\begin{proof}
显然由前面的引理 (i) 蕴含 (ii)，(ii) 蕴含 (iii)。

为证 (iii) 蕴含 (i)，对 $\sum_{j=1}^n\ell(V|_{U_j})$ 归纳。若此和为 $0$，则对所有 $j$ 有 $V|_{U_j}=0$，$V=0$。设此和严格为正。若 $V$ 不可约则已完成。否则存在非平凡拟凝聚 $\mathcal{D}_X$-子模 $U$，即有正合列
$$0\longrightarrow U\longrightarrow V\longrightarrow W\longrightarrow 0,$$
其中 $U$ 和 $W$ 均非零。由于对 $1\le j\le n$ 有 $\ell(V|_{U_j})=\ell(U|_{U_j})+\ell(W|_{U_j})$，可见
$$\sum_{j=1}^n\ell(V|_{U_j})=\sum_{j=1}^n\ell(U|_{U_j})+\sum_{j=1}^n\ell(W|_{U_j}),$$
且右边两项均不为零。故归纳假设可用于两者，由此可知 $U$ 和 $W$ 具有有限长度，从而 $V$ 具有有限长度。
\end{proof}

\section{特征簇}

本节把特征簇的构造推广到任意凝聚 $\mathcal{D}_X$-模。

先设 $X$ 是光滑仿射簇，$D_X$ 是 $X$ 上相应的微分算子环。由 II.3.7，$D_X$ 是左诺特且右诺特的环。此外它有按微分算子阶数的自然滤过 $F D_X$，由 II.3.6，作为配置这样滤过的滤过环满足 I.3 的公理。

设 $\pi:T^*(X)\to X$ 是 $X$ 的余切丛。由于 $\pi$ 是局部平凡纤维化，态射 $\pi$ 是仿射的，因此 $T^*(X)$ 是仿射簇。如 II.3.6 的证明中所述，有
$$\operatorname{Gr}D_X=\Gamma(X,\operatorname{Gr}\mathcal{D}_X)=\Gamma(X,\pi_*(\mathcal{O}_{T^*(X)}))=\Gamma(T^*(X),\mathcal{O}_{T^*(X)})=R(T^*(X)).$$

任意有限生成 $D_X$-模 $V_0$ 有好 $D_X$-模滤过 $F V_0$，且 $\Gr V_0$ 是 $\Gr D_X=R(T^*(X))$ 上的有限生成模。设 $I$ 是 $\Gr V_0$ 的零化子。则由 I.10.2，根 $r(I)$ 不依赖于 $V_0$ 上好滤过的选择。称之为 $V_0$ 的\emph{特征理想}，记为 $J(V_0)$。$J(V_0)$ 在 $T^*(X)$ 中的零点集称为 $V_0$ 的\emph{特征簇}，记为 $\operatorname{Ch}(V_0)$。这些定义与 I.10 中 $k^n$ 上微分算子模的定义一致。

现在把这些概念层化。设 $V$ 是 $X$ 上凝聚 $\mathcal{D}_X$-模，则称 $V$ 的特征簇 $\operatorname{Ch}(V)$ 为 $D_X$-模 $\Gamma(X,V)$ 的特征簇。

称 $V$ 由凝聚 $\mathcal{O}_X$-子模构成的递增 $\mathcal{D}_X$-模滤过 $F V$ 是\emph{好的}，若
\begin{enumerate}[label=(\roman*)]
  \item 对充分负的 $n\in\mathbb{Z}$，$F_nV=\{0\}$；
  \item 滤过 $F V$ 是穷尽的（即 $\bigcup_{n\in\mathbb{Z}}F_nV=V$）；
  \item 滤过 $F V$ 是稳定的，即存在 $m_0\in\mathbb{Z}$ 使得对所有 $n\in\mathbb{Z}_+$ 和 $m\ge m_0$ 有 $F_nD_X\,F_mV=F_{m+n}V$。
\end{enumerate}
设 $F V$ 是 $V$ 的好滤过，则 $\Gamma(X,F_pV)$ 是 $\Gamma(X,V)$ 的有限生成 $R(X)$-子模，且 $(\Gamma(X,F_pV);\ p\in\mathbb{Z})$ 是 $D_X$-模 $\Gamma(X,V)$ 的好滤过。

\begin{lemma}
设 $V$ 是凝聚 $\mathcal{D}_X$-模。则
\begin{enumerate}[label=(\roman*)]
  \item $V$ 有好滤过；
  \item 映射 $(F_pV;\ p\in\mathbb{Z})\longmapsto(\Gamma(X,F_pV);\ p\in\mathbb{Z})$ 是 $V$ 的好滤过之集到 $\Gamma(X,V)$ 的好滤过之集的双射。
\end{enumerate}
\end{lemma}

\begin{proof}
(i) $D_X$-模 $\Gamma(X,V)$ 有限生成，故它有好滤过 $F\Gamma(X,V)$。由 Serre 定理 $V=\Delta(\Gamma(X,V))$，$F_pV=\Delta(F_p\Gamma(X,V))$ 自然等同于 $V$ 的凝聚 $\mathcal{O}_X$-子模。直接检验可知 $F V$ 是 $V$ 上的好滤过。

(ii) 由 Serre 定理立即得出。
\end{proof}

设 $F V$ 是 $V$ 的好滤过，则 $\Gr V$ 是 $\Gr\mathcal{D}_X=\pi_*(\mathcal{O}_{T^*(X)})$ 上的模。此外由于 $X$ 仿射，有
$$\Gamma(X,\Gr V)=\Gr\Gamma(X,V),$$
其中 $\Gamma(X,V)$ 配置好滤过 $\{\Gamma(X,F_pV),\ p\in\mathbb{Z}\}$。由于 $\Gr\Gamma(X,V)$ 是有限生成 $R(T^*(X))$-模，由 I.4.2 对 $A=R(T^*(X))$ 的显然推广得
$$\operatorname{Ch}(V)=\operatorname{Ch}(\Gamma(X,V))=\operatorname{supp}(\Gr\Gamma(X,V))=\operatorname{supp}(\Gamma(X,\Gr V)).$$
这意味着特征簇的构造本质上是局部的，即有如下结果。设 $U$ 是 $X$ 中开集，则可把 $T^*(U)$ 看作 $T^*(X)$ 的开子集 $\pi^{-1}(U)=\{(x,\omega)\in T^*(X)\mid \omega\in T_x^*(X),\ x\in U\}$。

\begin{lemma}
设 $U$ 是 $X$ 中仿射开集。则
$$\operatorname{Ch}(V|_U)=\operatorname{Ch}(V)\cap\pi^{-1}(U).$$
\end{lemma}

现考虑一般情形。设 $X$ 是光滑代数簇，$\mathcal{D}_X$ 是 $X$ 上微分算子层，$\pi:T^*(X)\to X$ 是余切丛。设 $V$ 是凝聚 $\mathcal{D}_X$-模。

\begin{lemma}
存在唯一的 $T^*(X)$ 闭子簇 $M$，使得对任意仿射开集 $U\subset X$ 有 $\operatorname{Ch}(V|_U)=M\cap\pi^{-1}(U)$。
\end{lemma}

\begin{proof}
设 $U$ 和 $W$ 是 $X$ 的两个仿射开子集。则由前面的引理，
$$\operatorname{Ch}(V|_U)\cap\pi^{-1}(U\cap W)=\operatorname{Ch}(V|_{U\cap W})=\operatorname{Ch}(V|_W)\cap\pi^{-1}(U\cap W).$$
因此由所有满足如下条件的对 $(x,\omega)$（$\omega\in T_x^*(X)$）构成的集合 $M$ 是良定义的且具有所需性质：对 $x$ 的某仿射开邻域 $U$ 有 $(x,\omega)\in\operatorname{Ch}(V|_U)$。
\end{proof}

上一引理描述的簇 $M$ 即为 $V$ 的特征簇 $\operatorname{Ch}(V)$。

显然凝聚 $\mathcal{D}_X$-模好滤过的定义即使 $X$ 不仿射也有意义。现证明这种滤过的存在性。

先需要一个辅助结果。设 $X$ 是代数簇。证明凝聚 $\mathcal{O}_X$-子模可从开子簇扩张。

\begin{proposition}
设 $F$ 是 $X$ 上拟凝聚 $\mathcal{O}_X$-模，$G$ 是 $F|_U$ 的凝聚 $\mathcal{O}_U$-子模。则存在 $F$ 的凝聚 $\mathcal{O}_X$-子模 $G'$ 使得 $G'|_U=G$。
\end{proposition}

\begin{proof}
先设 $X$ 仿射。设 $\bar G$ 是 $G$ 的典范扩张。则 $\Gamma(X,\bar G)$ 是有限生成 $\Gamma(X,\mathcal{O}_X)$-模的递增族的直接极限。设 $\{H_i;\ i\in I\}$ 是这些子模的局部化。则它们构成 $F$ 的凝聚 $\mathcal{O}_X$-子模的递增系，其直接极限为 $\bar G$。由于 $\bar G|_U=G$ 凝聚，系 $\{H_i|_U;\ i\in I\}$ 稳定，即存在 $i_0\in I$ 使得对 $i\ge i_0$ 有 $H_i|_U=G$。

现考虑一般情形。对 $X$ 的有限仿射开覆盖的基数归纳证明。设 $(V_i;\ 1\le i\le n)$ 是 $X$ 的仿射开覆盖，$Y=\bigcup_{i=1}^{n-1}V_i$。则由归纳假设存在 $F|_Y$ 的凝聚 $\mathcal{O}_Y$-子模 $H$ 使得 $H|_{Y\cap U}=G|_{Y\cap U}$。$H$ 到限制于 $U$ 的 $F$ 的子模的典范扩张 $\bar H$ 含 $G$。应用证明的第一部分，存在 $\bar H|_{V_n}$ 的凝聚子模 $K$ 使得 $K|_{V_n\cap U}=G|_{V_n\cap U}$。设 $G'$ 为 $K$ 到 $\bar H$ 的子模的典范扩张。则 $G'|_U$ 含 $G$。此外 $G'|_{Y\cap U}\subset\bar H|_{Y\cap U}=G|_{Y\cap U}$，即 $G'|_{Y\cap U}=G|_{Y\cap U}$。又 $G'|_{V_n\cap U}=K|_{V_n\cap U}=G|_{V_n\cap U}$。因此 $G'|_U=G|_U$。另一方面，由于 $G'|_Y\subset\bar H|_Y=H$ 凝聚且 $G'|_{V_n}=K$ 也凝聚，$G'$ 是 $F$ 的凝聚 $\mathcal{O}_X$-子模。
\end{proof}

设 $X$ 是 $k$ 上的光滑代数簇，$\mathcal{D}_X$ 是 $X$ 上微分算子层。

\begin{theorem}
设 $V$ 是凝聚 $\mathcal{D}_X$-模。则 $V$ 有好滤过。
\end{theorem}

\begin{proof}
先断言存在凝聚 $\mathcal{O}_X$-子模 $U'\subset V$ 使得态射 $\mathcal{D}_X\otimes_{\mathcal{O}_X}U'\to V$ 是满射。设 $(U_i;\ 1\le i\le n)$ 是 $X$ 的仿射开覆盖。则对每个 $1\le i\le n$，$\Gamma(U_i,V)$ 是有限生成 $D_{U_i}$-模。由前面的命题，存在 $V$ 的凝聚 $\mathcal{O}_X$-子模 $G_i$ 使得 $\Gamma(U_i,G_i)$ 作为 $D_{U_i}$-模生成 $\Gamma(U_i,V)$。因此它们的和具有所需性质。

现可定义 $F_nV$ 为 $F_n\mathcal{D}_X\otimes_{\mathcal{O}_X}U'$ 在态射 $\mathcal{D}_X\otimes_{\mathcal{O}_X}U'\to V$ 下的像。显然 $F V$ 是 $V$ 的好滤过。
\end{proof}

设 $V$ 是凝聚 $\mathcal{D}_X$-模，$F V$ 是 $V$ 的好滤过。设 $\Gr V$ 是相应的分次 $\pi_*(\mathcal{O}_{T^*(X)})$-模。则对任意开集 $U\subset X$，滤过 $F_p(V|_U)=(F_pV)|_U$（$p\in\mathbb{Z}$）是 $V|_U$ 的好滤过。此外有 $\Gr V|_U=\Gr(V|_U)$。因此在 $X$ 的仿射开集 $U$ 上有 $\Gamma(U,\Gr V)=\Gr\Gamma(U,V)$。如前所述，簇 $T^*(U)$ 是仿射的且 $\Gamma(U,\pi_*(\mathcal{O}_{T^*(X)}))=R(T^*(U))$。故若把 $\Gamma(U,\Gr V)$ 作为 $R(T^*(U))$-模局部化，得到 $U$ 上唯一的 $\mathcal{O}_{T^*(U)}$-模 $\widetilde{V}_U$，具有性质 $\pi_*(\widetilde{V}_U)=\Gr V|_U$。由于 $\Gamma(U,\Gr V)$ 是有限生成 $R(T^*(U))$-模，$\widetilde{V}_U$ 是凝聚 $\mathcal{O}_{T^*(U)}$-模。把 $\widetilde{V}_U$ 粘在一起，得到 $T^*(X)$ 上唯一的凝聚 $\mathcal{O}_{T^*(X)}$-模 $\widetilde V$，具有性质 $\pi_*(\widetilde V)=\Gr V$。

这立即蕴含光滑仿射簇上凝聚模特征簇公式的如下推广。

\begin{proposition}
设 $V$ 是凝聚 $\mathcal{D}_X$-模，$F V$ 是 $V$ 的好滤过，$\Gr V$ 是相应的分次 $\pi_*(\mathcal{O}_{T^*(X)})$-模。则
\begin{enumerate}[label=(\roman*)]
  \item 存在唯一的 $T^*(X)$ 上凝聚 $\mathcal{O}_{T^*(X)}$-模 $\widetilde V$，使得 $\pi_*(\widetilde V)=\Gr V$；
  \item $\operatorname{Ch}(V)=\operatorname{supp}(\widetilde V)$。
\end{enumerate}
\end{proposition}

以下结果是 I.10.4 的推广。由于断言是局部的，可在仿射开集上检验，那里的论证与 I.10.4 相同。

\begin{proposition}
设
$$0\longrightarrow V_1\longrightarrow V_2\longrightarrow V_3\longrightarrow 0$$
是凝聚 $\mathcal{D}_X$-模的短正合列。则
$$\operatorname{Ch}(V_2)=\operatorname{Ch}(V_1)\cup\operatorname{Ch}(V_3).$$
\end{proposition}

设 $M$ 是 $X$ 的余切簇 $T^*(X)$ 的子簇。若 $(x,\omega)\in M$ 蕴含对所有 $\lambda\in k$ 有 $(x,\lambda\omega)\in M$，则称 $M$ 是\emph{锥形子簇}。

I.10.3 和 I.10.7 的推广由下列陈述给出。由于断言是局部的，只需在 II.2.10 意义下“小”的仿射开集 $U$ 上检验。此时 $\pi^{-1}(U)\cong U\times k^n$（$n=\dim_xX$），$R(T^*(U))\cong R(U)\otimes_kk[\xi_1,\dots,\xi_n]=R(U)[\xi_1,\dots,\xi_n]$，且分次是多项式环的自然分次。$\Gr\Gamma(U,V)$ 的零化子是 $R(U)[\xi_1,\dots,\xi_n]$ 中的齐次理想。在此情形下 I.10.3 和 I.10.7 证明的必要修改是直接的。

\begin{proposition}
设 $V$ 是 $X$ 上凝聚 $\mathcal{D}_X$-模。则特征簇 $\operatorname{Ch}(V)$ 是 $T^*(X)$ 的锥形子簇。
\end{proposition}

\begin{theorem}
设 $V$ 是 $X$ 上凝聚 $\mathcal{D}_X$-模。则
$$\pi(\operatorname{Ch}(V))=\operatorname{supp}(V).$$
\end{theorem}

\section{Coherentor（凝聚化函子）}

本节介绍代数簇上 $\mathcal{O}$-模的一些基本定义和结果。

先设 $X$ 是代数闭域 $k$ 上的任意代数簇。以 $\mathcal{M}(\mathcal{O}_X)$ 记 $X$ 上 $\mathcal{O}_X$-模的范畴，以 $\mathcal{M}_{qc}(\mathcal{O}_X)$ 记其拟凝聚 $\mathcal{O}_X$-模的全子范畴。整体截面函子 $\Gamma(X,-)$ 是从 $\mathcal{M}(\mathcal{O}_X)$ 到 $X$ 上正则函数环 $R(X)$ 上模范畴 $\mathcal{M}(R(X))$ 的左正合函子。由文献 [1] III.2.7，$\Gamma(X,-)$ 的右同调维数 $\le\dim X$。

\begin{lemma}
从 $\mathcal{M}_{qc}(\mathcal{O}_X)$ 到 $\mathcal{M}(\mathcal{O}_X)$ 的遗忘函子 $\operatorname{For}$ 有右伴随函子 $Q_X:\mathcal{M}(\mathcal{O}_X)\to\mathcal{M}_{qc}(\mathcal{O}_X)$。
\end{lemma}

\begin{proof}
只需证明对任意 $W\in\mathcal{M}(\mathcal{O}_X)$ 存在 $Q(W)\in\mathcal{M}_{qc}(\mathcal{O}_X)$ 使得
$$\operatorname{Hom}_{\mathcal{O}_X}(V,W)=\operatorname{Hom}_{\mathcal{O}_X}(V,Q(W))$$
对任意 $V\in\mathcal{M}_{qc}(\mathcal{O}_X)$ 成立。

先设 $X$ 是仿射簇。则 $\Gamma(X,-)$ 是 $\mathcal{M}_{qc}(\mathcal{O}_X)$ 与 $\mathcal{M}(R(X))$ 的范畴等价。此外对任意 $R(X)$-模 $M$，以 $\widetilde M=\mathcal{O}_X\otimes_{R(X)}M$ 记其局部化。则 $M\to\widetilde M$ 是从 $\mathcal{M}(R(X))$ 到 $\mathcal{M}_{qc}(\mathcal{O}_X)$ 的正合函子，且是 $\Gamma(X,-)$ 的拟逆。因此对任意 $V\in\mathcal{M}_{qc}(\mathcal{O}_X)$，
\begin{align*}
\operatorname{Hom}_{\mathcal{O}_X}(V,W)&=\operatorname{Hom}_{\mathcal{O}_X}(\widetilde{\Gamma(X,V)},W)\\
&=\operatorname{Hom}_{\mathcal{O}_X}(\mathcal{O}_X\otimes_{R(X)}\Gamma(X,V),W)=\operatorname{Hom}_{R(X)}(\Gamma(X,V),\operatorname{Hom}_{\mathcal{O}_X}(\mathcal{O}_X,W))\\
&=\operatorname{Hom}_{R(X)}(\Gamma(X,V),\Gamma(X,W))=\operatorname{Hom}_{\mathcal{O}_X}(V,\widetilde{\Gamma(X,W)}),
\end{align*}
此时 $Q_X(W)=\widetilde{\Gamma(X,W)}$。

现设 $U$ 是仿射簇 $X$ 的开仿射子集，$i:U\to X$ 是自然浸入。则对任意 $V\in\mathcal{M}_{qc}(\mathcal{O}_X)$ 和 $W\in\mathcal{M}(\mathcal{O}_U)$，
$$\operatorname{Hom}_{\mathcal{O}_X}(V,i_*(W))=\operatorname{Hom}_{\mathcal{O}_U}(V|_U,W)=\operatorname{Hom}_{\mathcal{O}_U}(V|_U,Q_U(W))=\operatorname{Hom}_{\mathcal{O}_X}(V,i_*(Q_U(W))).$$
由于直像保持拟凝聚性（文献 [1] II.5.8），$i_*(Q_U(W))\in\mathcal{M}_{qc}(\mathcal{O}_X)$，且
$$Q_X(i_*(W))=i_*(Q_U(W)),$$
即函子 $Q_X\circ i_*$ 与 $i_*\circ Q_U$ 同构。

现考虑一般情形。设 $X$ 是任意簇，$U=(U_1,U_2,\dots,U_n)$ 是其仿射开覆盖。以 $f_i:U_i\to X$，$f_{ij}:U_i\cap U_j\to X$ 记自然浸入。由相应的 Čech 分解，对任意 $W\in\mathcal{M}(\mathcal{O}_X)$ 有正合列
$$0\longrightarrow W\longrightarrow\bigoplus_{i=1}^n f_{i*}(W|_{U_i})\longrightarrow\bigoplus_{i<j}f_{ij*}(W|_{U_i\cap U_j}).$$
固定 $1\le i\le n$，以 $g_{ij}:U_i\cap U_j\to U_i$ 记自然包含（$j\neq i$）。则 $f_{ij}=f_i\circ g_{ij}$，且 $f_{i*}(W|_{U_i})\to\bigoplus_{i\neq j}f_{ij*}(W|_{U_i\cap U_j})$ 是把 $f_i$ 的直像应用于 $W|_{U_i}\to\bigoplus_{j\neq i}g_{ij*}(W|_{U_i\cap U_j})$ 得到的。由于 $U_i$ 和 $U_i\cap U_j$ 是开仿射子簇，对此态射应用 $Q_{U_i}$ 并用前段结果，得到态射 $Q_{U_i}(W|_{U_i})\to\bigoplus_{j\neq i}g_{ij*}(Q_{U_i\cap U_j}(W|_{U_i\cap U_j}))$，使得对任意 $V\in\mathcal{M}_{qc}(\mathcal{O}_X)$ 图
$$\begin{tikzcd}
\operatorname{Hom}_{\mathcal{O}_X}(V,f_{i*}(W|_{U_i})) \arrow[r] \arrow[d] & \operatorname{Hom}_{\mathcal{O}_X}(V,f_{ij*}(W|_{U_i\cap U_j})) \arrow[d]\\
\operatorname{Hom}_{\mathcal{O}_{U_i}}(V|_{U_i},W|_{U_i}) \arrow[r] \arrow[d] & \operatorname{Hom}_{\mathcal{O}_{U_i}}(V|_{U_i},g_{ij*}(W|_{U_i\cap U_j})) \arrow[d]\\
\operatorname{Hom}_{\mathcal{O}_{U_i}}(V|_{U_i},Q_{U_i}(W|_{U_i})) \arrow[r] & \operatorname{Hom}_{\mathcal{O}_{U_i}}(V|_{U_i},g_{ij*}(Q_{U_i\cap U_j}(W|_{U_i\cap U_j})))\\
\operatorname{Hom}_{\mathcal{O}_X}(V,f_{i*}(Q_{U_i}(W|_{U_i}))) \arrow[u] \arrow[r] & \operatorname{Hom}_{\mathcal{O}_X}(V,f_{ij*}(Q_{U_i\cap U_j}(W|_{U_i\cap U_j}))) \arrow[u]
\end{tikzcd}$$
交换。Čech 分解 $C^\cdot(U,W)$ 的第一个微分
$$d^1:\bigoplus_{i=1}^nf_{i*}(W|_{U_i})\longrightarrow\bigoplus_{i<j}f_{ij*}(W|_{U_i\cap U_j})$$
由此对应决定态射
$$\delta:\bigoplus_{i=1}^nf_{i*}(Q_{U_i}(W|_{U_i}))\longrightarrow\bigoplus_{i<j}f_{ij*}(Q_{U_i\cap U_j}(W|_{U_i\cap U_j})).$$
以 $Q(W)$ 记此态射的核。显然它是拟凝聚 $\mathcal{O}_X$-模。由函子 $\operatorname{Hom}_{\mathcal{O}_X}(V,-)$ 的左正合性得 $\operatorname{Hom}_{\mathcal{O}_X}(V,W)=\operatorname{Hom}_{\mathcal{O}_X}(V,Q(W))$。
\end{proof}

函子 $Q_X:\mathcal{M}(\mathcal{O}_X)\to\mathcal{M}_{qc}(\mathcal{O}_X)$ 称为\emph{凝聚化函子}（coherentor）。

\begin{proposition}
\begin{enumerate}[label=(\roman*)]
  \item 函子 $Q_X:\mathcal{M}(\mathcal{O}_X)\to\mathcal{M}_{qc}(\mathcal{O}_X)$ 是左正合的；
  \item $Q_X$ 把 $\mathcal{M}(\mathcal{O}_X)$ 中的入射对象映为 $\mathcal{M}_{qc}(\mathcal{O}_X)$ 中的入射对象；
  \item 复合 $Q_X\circ\operatorname{For}$ 同构于 $\mathcal{M}_{qc}(\mathcal{O}_X)$ 上的恒等函子；
  \item 设 $U$ 是 $X$ 的开子簇，$i:U\to X$ 是自然包含。则函子 $Q_X\circ i_*$ 与 $i_*\circ Q_U$ 同构。
\end{enumerate}
\end{proposition}

\begin{proof}
(i) 这是任意右伴随函子的性质。

(ii) 这是正合函子的右伴随的性质。

(iii) 由 $Q_X$ 的定义显然。

(iv) 对任意 $V\in\mathcal{M}_{qc}(\mathcal{O}_X)$ 和 $W\in\mathcal{M}(\mathcal{O}_U)$ 有
$$\operatorname{Hom}_{\mathcal{O}_X}(V,i_*(W))=\operatorname{Hom}_{\mathcal{O}_U}(V|_U,W)=\operatorname{Hom}_{\mathcal{O}_U}(V|_U,Q_U(W))=\operatorname{Hom}_{\mathcal{O}_X}(V,i_*(Q_U(W))).$$
由于直像保持拟凝聚性，$i_*(Q_U(W))\in\mathcal{M}_{qc}(\mathcal{O}_X)$，且 $Q_X(i_*(W))=i_*(Q_U(W))$，即函子 $Q_X\circ i_*$ 与 $i_*\circ Q_U$ 同构。
\end{proof}

\begin{theorem}
范畴 $\mathcal{M}_{qc}(\mathcal{O}_X)$ 有充足的入射对象。
\end{theorem}

\begin{proof}
设 $V$ 是拟凝聚 $\mathcal{O}_X$-模。则存在入射 $\mathcal{O}_X$-模 $I$ 和单射 $\epsilon:V\to I$。由于 $Q_X$ 由前面的命题 (i) 左正合，$Q_X(\epsilon):Q_X(V)\to Q_X(I)$ 是单射。此外由 (ii)，$Q_X(I)$ 在 $\mathcal{M}_{qc}(\mathcal{O}_X)$ 中入射。最后由 (iii)，$Q_X(V)\cong V$。
\end{proof}

\begin{lemma}
设 $X$ 是仿射簇，$W\in\mathcal{M}(\mathcal{O}_X)$。则 $R^pQ_X(W)=\widetilde{H^p(X,W)}$。
\end{lemma}

\begin{proof}
在引理 4.1 的证明中已见 $Q_X$ 同构于局部化函子与整体截面函子 $\Gamma(X,-)$ 的复合。设 $I^\cdot$ 是 $W$ 的入射分解。则 $H^p(X,W)=H^p(\Gamma(X,I^\cdot))$。由于局部化函子正合，
$$\widetilde{H^p(X,W)}=H^p(\widetilde{\Gamma(X,I^\cdot)})=H^p(Q_X(I^\cdot))=R^pQ_X(W).$$
\end{proof}

以下引理在凝聚化函子各种性质的证明中至关重要。

\begin{lemma}
设 $V$ 是满足如下条件的 $\mathcal{O}_X$-模：

(V) 对任意仿射开集 $U\subset X$ 和 $p\ge 1$ 有上同调 $H^p(U,V)=0$。

则对任意仿射开集 $U\subset X$，直像层 $i_*(V|_U)$ 是 $Q_X$-无环的。
\end{lemma}

\begin{proof}
先注意高阶直像 $R^pi_*(V|_U)$ 为零。固定 $p>0$。对 $U$ 上任意层 $F$，$R^pi_*(F)$ 是预层 $V\longmapsto H^p(i^{-1}(V),F)=H^p(U\cap V,F)$ 附着的层。若 $V$ 是仿射开集，$U\cap V$ 也是仿射的，故 $H^p(U\cap V,V)=0$。这意味着此预层在所有仿射开集上为零。由于仿射开集构成 $X$ 拓扑的基，相应的层为零。

设 $I^\cdot$ 是 $V|_U$ 在 $\mathcal{M}(\mathcal{O}_X)$ 中的入射分解。由前注复形 $i_*(I^\cdot)$ 是 $i_*(V|_U)$ 的分解。此外由于 $i_*$ 是限制函子 $\mathcal{M}(\mathcal{O}_X)\to\mathcal{M}(\mathcal{O}_U)$ 的右伴随，它把入射映为入射，即 $i_*(I^\cdot)$ 是 $i_*(V|_U)$ 在 $\mathcal{M}(\mathcal{O}_X)$ 中的入射分解。因此
$$R^pQ_X(i_*(V|_U))=H^p(Q_X(i_*(I^\cdot)))=H^p(i_*(Q_U(I^\cdot))),$$
由命题 (iv)。由 4.4 对 $p>0$，
$$H^p(Q_U(I^\cdot))=R^pQ_U(V|_U)=\widetilde{H^p(U,V)}=0.$$
因此 $Q_U(I^\cdot)$ 是由拟凝聚 $\mathcal{O}_U$-模构成的无环复形。由于 $i:U\to X$ 是仿射态射，$i_*$ 是从 $\mathcal{M}_{qc}(\mathcal{O}_U)$ 到 $\mathcal{M}_{qc}(\mathcal{O}_X)$ 的正合函子。故复形 $i_*(Q_U(I^\cdot))$ 也无环，即 $R^pQ_X(i_*(V|_U))=0$（$p>0$）。
\end{proof}

\begin{proposition}
拟凝聚 $\mathcal{O}_X$-模对 $Q_X$ 无环。
\end{proposition}

\begin{proof}
设 $U=(U_0,U_1,\dots,U_n)$ 是 $X$ 的仿射开覆盖。设 $V$ 是拟凝聚 $\mathcal{O}_X$-模，则 $V$ 满足前一引理的条件 (V)。因此 $V$ 的 Čech 复形 $C^\cdot(U,V)$ 中的模都是 $Q_X$-无环的。可知 $C^\cdot(U,V)$ 是 $V$ 的 $Q_X$-无环分解。故由命题 (iii)，
$$R^pQ_X(V)=H^p(Q_X(C^\cdot(U,V)))=H^p(C^\cdot(U,V))=0\quad(p>0).$$
\end{proof}

此结果蕴含如下基本事实。

\begin{theorem}
设 $V,W$ 是两个拟凝聚 $\mathcal{O}_X$-模。则
$$\operatorname{Ext}^\cdot_{\mathcal{M}_{qc}(\mathcal{O}_X)}(V,W)=\operatorname{Ext}^\cdot_{\mathcal{M}(\mathcal{O}_X)}(V,W).$$
\end{theorem}

\begin{proof}
设
$$0\longrightarrow W\longrightarrow I^0\longrightarrow I^1\longrightarrow\cdots\longrightarrow I^n\longrightarrow\cdots$$
是 $W$ 在 $\mathcal{M}(\mathcal{O}_X)$ 中的入射分解。则由命题 (ii)、(iii) 和前一命题，
$$0\longrightarrow W\longrightarrow Q_X(I^0)\longrightarrow Q_X(I^1)\longrightarrow\cdots\longrightarrow Q_X(I^n)\longrightarrow\cdots$$
是 $W$ 在 $\mathcal{M}_{qc}(\mathcal{O}_X)$ 中的入射分解。故对任意 $p\in\mathbb{Z}_+$，
$$\operatorname{Ext}^p_{\mathcal{M}(\mathcal{O}_X)}(V,W)=H^p(\operatorname{Hom}_{\mathcal{O}_X}(V,I^\cdot))=H^p(\operatorname{Hom}_{\mathcal{O}_X}(V,Q_X(I^\cdot)))=\operatorname{Ext}^p_{\mathcal{M}_{qc}(\mathcal{O}_X)}(V,W).$$
\end{proof}

因此可无混淆地记
$$\operatorname{Ext}^\cdot_{\mathcal{O}_X}(V,W)=\operatorname{Ext}^\cdot_{\mathcal{M}_{qc}(\mathcal{O}_X)}(V,W)=\operatorname{Ext}^\cdot_{\mathcal{M}(\mathcal{O}_X)}(V,W)$$
对任意两个拟凝聚 $\mathcal{O}_X$-模 $V,W$。

\begin{theorem}
$Q_X$ 的右同调维数有限。
\end{theorem}

\begin{proof}
固定 $X$ 的仿射开覆盖 $U=(U_0,U_1,\dots,U_n)$。

先观察如下事实：设 $V$ 是满足引理 4.5 条件 (V) 的 $\mathcal{O}_X$-模，则对 $p>n$ 有 $R^pQ_X(V)=0$。

如前证明，由引理 4.5，$V$ 的 Čech 复形 $C^\cdot(U,V)$ 中的模都是 $Q_X$-无环的，故 $C^\cdot(U,V)$ 是 $Q_X$-无环分解。因此 $R^pQ_X(V)=H^p(Q_X(C^\cdot(U,V)))$，由此对 $p>n$ 有 $R^pQ_X(V)=0$。

现建立此结果的如下推广：设 $V$ 是满足如下条件的 $\mathcal{O}_X$-模：

(V$_q$) 对任意仿射开集 $U\subset X$ 和 $p>q$ 有 $H^p(U,V)=0$。

则对 $p>q+n$ 有 $R^pQ_X(V)=0$。

已对 $q=0$ 建立此结果。为证归纳步，对 $p$ 归纳。设结论对某 $q\ge 0$ 成立。设 $V$ 是使 $H^p(U,V)=0$（对 $p>q+1$ 和任意仿射开集 $U\subset X$）的 $\mathcal{O}_X$-模。取短正合列
$$0\longrightarrow V\longrightarrow I\longrightarrow W\longrightarrow 0,$$
其中 $I$ 是入射 $\mathcal{O}_X$-模。则对任意仿射开集 $U\subset X$，限制 $I|_U$ 是入射 $\mathcal{O}_U$-模。因此由上同调长正合列得连接同态 $H^p(U,W)\to H^{p+1}(U,V)$ 对 $p\ge 1$ 是同构。这意味着 $H^p(U,W)=0$（$p>q$），即归纳假设可用于 $W$。由 $Q_X$ 的导出函子的长正合列得连接态射 $R^pQ_X(W)\to R^{p+1}Q_X(V)$ 对 $p\ge 1$ 是同构。由于 $p>q+n$ 时 $R^pQ_X(W)=0$，$p>q+1+n$ 时 $R^pQ_X(V)=0$。故 (V$_{q+1}$) 成立。

由于函子 $\Gamma(U,-)$ 的右同调维数 $\le\dim U\le\dim X$，可见所有 $\mathcal{O}_X$-模满足 (V$_{\dim X}$)。因此对任意拟凝聚 $\mathcal{O}_X$-模 $V$，当 $p>\dim X+n$ 时 $R^pQ_X(V)=0$。
\end{proof}
\begin{annote}
$\mathcal Q$（coherator）是 $\mathcal D_X$-模层到凝聚 $\mathcal D_X$-模的右伴随，$\mathcal Q(\mathcal M)=\varinjlim\mathcal H om_{\mathcal D_X}(\mathcal E,\mathcal M)$，$\mathcal E$ 取遍凝聚 $\mathcal D_X$-模。它修补了“$\mathcal D_X$ 非凝聚导致的层论不灵”，使 $\mathcal D_X$-模可以直接在 $D^+(\mathcal D_X)$ 里做内射消解。Stacks 与 nLab 均无独立条目，本条只作原书内部的记号澄清。[来源：原书 §4, p90–96；译者按]
\end{annote}



\section{射影空间上的 $D$-模}

设 $X=P^n$ 是 $k$ 上 $n$ 维射影空间。记 $Y=k^{n+1}$，$Y^*=Y-\{0\}$。设 $\pi:Y^*\to X$ 是自然投射
$$\pi(x_0,x_1,\dots,x_n)=[x_0,x_1,\dots,x_n].$$
令 $U_0=\{[x_0,x_1,\dots,x_n]\in P^n\mid x_0\neq 0\}$。则 $U_0$ 是 $P^n$ 中同构于 $k^n$ 的开集，同构 $c:k^n\to U_0$ 由 $c(x_1,x_2,\dots,x_n)=[1,x_1,x_2,\dots,x_n]$ 给出。此外 $\pi^{-1}(U_0)=\{(y_0,y_1,y_2,\dots,y_n)\mid y_0\neq 0\}$。因此由
$$\varphi(t,z_1,z_2,\dots,z_n)=(t,tz_1,tz_2,\dots,tz_n)$$
定义的映射 $\varphi:k^*\times k^n\to\pi^{-1}(U_0)$ 是 $k^*\times k^n$ 与 $\pi^{-1}(U_0)$ 的同构，且满足 $\pi(\varphi(t,z_1,\dots,z_n))=[1,z_1,\dots,z_n]$。因此如下交换图交换：
$$\begin{tikzcd}
k^*\times k^n \arrow[r,"\varphi"] \arrow[d,"1\times c"' ] & \pi^{-1}(U_0) \arrow[d,"\pi"]\\
k^*\times U_0 \arrow[r,"\operatorname{pr}_2"'] & U_0
\end{tikzcd}$$
显然 $\alpha=\varphi\circ(1\times c)^{-1}:k^*\times U_0\to\pi^{-1}(U_0)$ 是由
$$\alpha(t,[1,x_1,x_2,\dots,x_n])=(t,tx_1,tx_2,\dots,tx_n)$$
给出的同构，它在开集 $U_0$ 上平凡化纤维化 $\pi:Y^*\to X$。由于 $\operatorname{GL}(n+1,k)$ 在 $P^n$ 上传递作用，这证明如下断言。

\begin{lemma}
态射 $\pi:Y^*\to X$ 是纤维同构于 $k^*$ 的局部平凡纤维化。
\end{lemma}

显然 $U_0$ 和 $\pi^{-1}(U_0)$ 是仿射簇。因此若 $V$ 是拟凝聚 $\mathcal{D}_X$-模，有
$$\Gamma\bigl(\pi^{-1}(U_0),\pi^+(V)\bigr)=R\bigl(\pi^{-1}(U_0)\bigr)\otimes_{R(U_0)}\Gamma(U_0,V).$$
另一方面，同构 $\alpha:k^*\times U_0\to\pi^{-1}(U_0)$ 诱导同构 $\alpha^*:R(\pi^{-1}(U_0))\to R(k^*)\otimes_kR(U_0)$。在此同构下，$R(U_0)$-模作用对应于第二个因子中的乘法。因此此同构诱导 $\Gamma(\pi^{-1}(U_0),\pi^+(V))$ 与 $R(k^*)\otimes_k\Gamma(U_0,V)$ 的同构。

再次利用 $\operatorname{GL}(n+1,k)$ 作用的传递性，立即得到如下推论。

\begin{lemma}
逆像函子 $\pi^+:\mathcal{M}_{qc}(\mathcal{D}_X)\to\mathcal{M}_{qc}(\mathcal{D}_{Y^*})$ 是正合的。
\end{lemma}

因此若
$$0\longrightarrow V_1\longrightarrow V_2\longrightarrow V_3\longrightarrow 0$$
是拟凝聚 $\mathcal{D}_X$-模的正合列，有拟凝聚 $\mathcal{D}_{Y^*}$-模的正合列
$$0\longrightarrow\pi^+(V_1)\longrightarrow\pi^+(V_2)\longrightarrow\pi^+(V_3)\longrightarrow 0.$$
此外若以 $j:Y^*\to Y$ 记自然浸入，有拟凝聚 $\mathcal{D}_Y$-模的正合列
$$0\longrightarrow j_+(\pi^+(V_1))\longrightarrow j_+(\pi^+(V_2))\longrightarrow j_+(\pi^+(V_3))\longrightarrow K\longrightarrow 0,$$
且 $K$ 支集在 $\{0\}$。由于 $Y$ 是仿射空间，这意味着
$$0\to\Gamma(Y,j_+(\pi^+(V_1)))\to\Gamma(Y,j_+(\pi^+(V_2)))\to\Gamma(Y,j_+(\pi^+(V_3)))\to\Gamma(Y,K)\to 0$$
是 $D(n+1)$-模的正合列。对任意 $\mathcal{D}_{Y^*}$-模 $W$，由于 $j_+(W)=j_\cdot(W)$，有 $\Gamma(Y,j_+(W))=\Gamma(Y^*,W)$。故断定
$$0\to\Gamma(Y^*,\pi^+(V_1))\to\Gamma(Y^*,\pi^+(V_2))\to\Gamma(Y^*,\pi^+(V_3))\to\Gamma(Y,K)\to 0$$
是 $D(n+1)$-模的正合列。

令 $E=\sum_{i=0}^ny_i\dfrac{\partial}{\partial y_i}$ 为 $Y$ 上 Euler 算子。微分算子 $E$ 是 $Y$ 上向量场。若 $y=(y_0,y_1,\dots,y_n)$ 是 $Y^*$ 中一点，$E(y)$ 是曲线 $t\longmapsto(ty_0,ty_1,\dots,ty_n)$ 在 $y$ 处的切向量。故有如下结果。

\begin{lemma}
对任意 $y\in Y^*$，Euler 算子 $E$ 在 $y$ 处的值属于微分 $\pi:Y^*\to X$ 的核。
\end{lemma}

在同构 $\alpha:k^*\times U_0\to\pi^{-1}(U_0)$ 下，Euler 算子对应于 $k^*\times U_0$ 上的微分算子 $t\partial_t$，其中 $t$ 是 $k^*$ 上自然坐标。显然 $R(k^*)$ 是环 $k((t))$，即多项式环 $k[t]$ 关于乘法系 $t^n$（$n\in\mathbb{Z}_+$）的局部化。因此在 $\alpha$ 给出的同构下，
$$\Gamma\bigl(\pi^{-1}(U_0),\pi^+(V)\bigr)\cong k((t))\otimes_k\Gamma(U_0,V).$$
因此 $\Gamma(\pi^{-1}(U_0),\pi^+(V))$ 的每个截面被 $\prod_{p\in S}(E-p)$ 零化，其中 $S$ 是依赖于该截面的 $\mathbb{Z}$ 的有限子集。

\begin{lemma}
设 $V$ 是拟凝聚 $\mathcal{D}_X$-模。则 $\Gamma(Y^*,\pi^+(V))$ 是特征值属于 $\mathbb{Z}$ 的 $E$-特征空间的直和。
\end{lemma}

\begin{proof}
由于 $U_0$ 在 $\operatorname{GL}(n+1,k)$ 作用下的有限多个平移覆盖 $X$，$\pi^{-1}(U_0)$ 的有限多个平移覆盖 $Y^*$。由上述论证，对 $V$ 的任意整体截面 $v$，存在有限子集 $S\subset\mathbb{Z}$ 使 $\prod_{p\in S}(E-p)v=0$。
\end{proof}

因此所考虑的正合列在 $E$ 作用下分裂为对应于 $E$ 不同特征值的无穷族正合列。特别地若以 $\Gamma(Y^*,\pi^+(V))_{(p)}$ 记 $\Gamma(Y^*,\pi^+(V))$ 中特征值 $p\in\mathbb{Z}$ 的 $E$-特征空间，则对任意 $p\in\mathbb{Z}$，
$$0\to\Gamma(Y^*,\pi^+(V_1))_{(p)}\to\Gamma(Y^*,\pi^+(V_2))_{(p)}\to\Gamma(Y^*,\pi^+(V_3))_{(p)}\to\Gamma(Y,K)_{(p)}\to 0.$$

另一方面，有自然线性映射 $\gamma_U:\Gamma(U,V)\to\Gamma(\pi^{-1}(U),\pi^+(V))$，它与限制相容，即对 $X$ 中任意两个开集 $V\subset U$ 图
$$\begin{tikzcd}
\Gamma(U,V) \arrow[r,"\gamma_U"] \arrow[d,"r_{V,U}"'] & \Gamma(\pi^{-1}(U),\pi^+(V)) \arrow[d,"r_{\pi^+(V),\pi^{-1}(U)}"]\\
\Gamma(V,V) \arrow[r,"\gamma_V"'] & \Gamma(\pi^{-1}(V),\pi^+(V))
\end{tikzcd}$$
交换。显然在上述等同下 $\gamma_{U_0}$ 对应于映射 $v\longmapsto 1\otimes v|_{U_0}$（从 $\Gamma(U_0,V)$ 到 $k((t))\otimes_k\Gamma(U_0,V)$）。因此它是 $\Gamma(U_0,V)$ 到 $\Gamma(\pi^{-1}(U_0),\pi^+(V))_{(0)}$ 的同构。故可把 $\gamma=\gamma_X$ 看作从 $\Gamma(X,V)$ 到 $\Gamma(Y^*,\pi^+(V))_{(0)}$ 的线性映射。

显然 $\gamma(v)=0$ 蕴含 $\gamma(v)|_{U_0}=0$，从而 $v|_{U_0}=0$。故再次利用 $\operatorname{GL}(n+1,k)$ 在 $X$ 上作用的传递性得 $v=0$。由此可知 $\gamma:\Gamma(X,V)\to\Gamma(Y^*,\pi^+(V))_{(0)}$ 是单射。

\begin{lemma}
设 $V$ 是拟凝聚 $\mathcal{D}_X$-模。则 $\gamma:\Gamma(X,V)\to\Gamma(Y^*,\pi^+(V))_{(0)}$ 是同构。
\end{lemma}

\begin{proof}
设 $s$ 是 $\Gamma(Y^*,\pi^+(V))_{(0)}$ 的整体截面。则由上述讨论，其在 $\pi^{-1}(U_0)$ 上的限制等于某截面 $v_0\in V(U_0)$ 的 $\gamma_{U_0}(v_0)$。由于 $U_0$ 在 $\operatorname{GL}(n+1,k)$ 作用下的有限多个平移覆盖 $X$，可见存在 $X$ 的覆盖 $U_0,\dots,U_m$ 和 $V$ 在这些开集上的截面 $v_0,\dots,v_m$，使得 $\gamma_{U_k}(v_k)=s|_{\pi^{-1}(U_k)}$（$0\le k\le m$）。由于 $\gamma_{U_k}$ 是单射且 $U_k$ 和 $\pi^{-1}(U_k)$ 是仿射簇，由局部化得 $\gamma_{U_k\cap U_l}$ 是单射。故对每对 $0\le k<l\le m$ 有 $v_k|_{U_k\cap U_l}=v_l|_{U_k\cap U_l}$。由此可知存在 $V$ 的整体截面 $v$ 使得对所有 $0\le k\le m$ 有 $v|_{U_k}=v_k$。因此
$$\gamma(v)|_{\pi^{-1}(U_k)}=\gamma_{U_k}(v|_{U_k})=\gamma_{U_k}(v_k)=s|_{\pi^{-1}(U_k)}$$
对 $0\le k\le m$ 成立，$\gamma(v)=s$。这证明 $\gamma$ 的满射性。
\end{proof}

以下结果由直接计算得出。

\begin{lemma}
对任意 $0\le i\le n$，
\begin{enumerate}[label=(\roman*)]
  \item $[E,y_i]=y_i$；
  \item $\bigl[E,\dfrac{\partial}{\partial y_i}\bigr]=-\dfrac{\partial}{\partial y_i}$。
\end{enumerate}
\end{lemma}

\begin{lemma}
设 $Y=k^{n+1}$，$V$ 是支集在 $\{0\}$ 的拟凝聚 $\mathcal{D}_Y$-模。则 $\Gamma(Y,V)$ 是特征值属于 $\{-(n+k)\mid k\in\mathbb{N}\}$ 的 $E$-特征空间的直和。
\end{lemma}

\begin{proof}
设 $V_0=\Gamma(Y,V)$。由 I.13.7 知 $V_0$ 由子空间 $W$ 生成，其中 $W$ 由被 $y_i$（$0\le i\le n$）零化的所有整体截面构成。取 $v\in W$，则 $v=[\partial_i,y_i]v=-y_i\partial_iv$，故 $Ev=-(n+1)v$。此外 $V_0=\bigoplus_{I\in\mathbb{Z}_+^{n+1}}\partial^I(W)$。由前面的结果，对任意 $w=\partial^Iv$（$v\in W$）有
$$Ew=E\partial^Iv=\partial^IEv-|I|\partial^Iv=\bigl(-(n+1)-|I|\bigr)\partial^Iv=-(n+1+|I|)w,$$
即 $V_0$ 是 $E$ 的属于特征值 $\{-(n+k)\mid k\in\mathbb{N}\}$ 的特征空间的直和。
\end{proof}

特别地，由于 $K$ 支集在 $\{0\}$，有 $\Gamma(Y,K)_{(0)}=\{0\}$，且
$$0\longrightarrow\Gamma(X,V_1)\longrightarrow\Gamma(X,V_2)\longrightarrow\Gamma(X,V_3)\longrightarrow 0$$
正合。因此 $\Gamma(X,-)$ 在 $\mathcal{M}_{qc}(\mathcal{D}_X)$ 上是正合函子。

\begin{theorem}
设 $V$ 是 $n$ 维射影空间 $X=P^n$ 上拟凝聚 $\mathcal{D}_X$-模。则对 $p\in\mathbb{N}$ 有 $H^p(X,V)=0$。
\end{theorem}

\begin{proof}
$V$ 有由入射拟凝聚 $\mathcal{D}_X$-模构成的右分解 $I^\cdot$。入射拟凝聚 $\mathcal{D}_X$-模是 $\Gamma$-无环的。因此 $H^p(X,V)=H^p(\Gamma(X,I^\cdot))=0$（$p>0$），因为 $\Gamma$ 在拟凝聚 $\mathcal{D}_X$-模上正合。
\end{proof}

因此若以 $D_X=\Gamma(X,\mathcal{D}_X)$ 记 $X$ 上整体微分算子，$\Gamma:\mathcal{M}_{qc}(\mathcal{D}_X)\to\mathcal{M}(D_X)$ 是正合函子。还可定义函子 $\Delta:\mathcal{M}(D_X)\to\mathcal{M}(D_X)$ 为 $\Delta(V)=D_X\otimes_{D_X}V$，则 $\Delta$ 是右正合的。设 $V\in\mathcal{M}(D_X)$，则存在正合列 $D_X^{(I)}\to D_X^{(J)}\to V\to 0$，应用 $\Delta$ 后得到 $D_X$-模的正合列 $D_X^{(I)}\to D_X^{(J)}\to\Delta(V)\to 0$。因此 $\Delta(V)$ 是拟凝聚 $\mathcal{D}_X$-模。

函子 $\Gamma\circ\Delta$ 是 $\mathcal{M}(D_X)$ 到自身的右正合函子。此外对任意 $V\in\mathcal{M}(D_X)$ 存在自然态射 $\lambda_V:V\to\Gamma(X,\Delta(V))$。显然 $\lambda$ 是恒等函子到 $\Gamma\circ\Delta$ 的自然变换。

\begin{lemma}
自然变换 $\lambda$ 是 $\mathcal{M}(D_X)$ 上恒等函子到函子 $\Gamma\circ\Delta$ 的同构。
\end{lemma}

\begin{proof}
显然对任意自由 $D_X$-模 $F$，$\lambda_F:F\to\Gamma(X,\Delta(F))$ 是同构。因此取 $D_X$-模的正合列 $D_X^{(I)}\to D_X^{(J)}\to V\to 0$，得到如下 $D_X$-模的交换图
$$\begin{tikzcd}
D_X^{(I)} \arrow[r] \arrow[d,"\lambda_{D_X^{(I)}}"'] & D_X^{(J)} \arrow[r] \arrow[d,"\lambda_{D_X^{(J)}}"'] & V \arrow[r] \arrow[d,"\lambda_V"] & 0\\
D_X^{(I)} \arrow[r] & D_X^{(J)} \arrow[r] & \Gamma(X,\Delta(V)) \arrow[r] & 0
\end{tikzcd}$$
其行正合且前两个竖直箭头是同构，因此 $\lambda_V$ 是同构。
\end{proof}

\begin{lemma}
设 $V$ 是拟凝聚 $\mathcal{D}_X$-模。若 $\Gamma(X,V)=0$，则 $V=0$。
\end{lemma}

\begin{proof}
设 $\Gamma(X,V)=0$。由前面的结果这意味着 $\Gamma(Y^*,\pi^+(V))_{(0)}=0$。断言实际上 $\Gamma(Y^*,\pi^+(V))=0$。反设不然，则存在 $s\in\mathbb{Z}-\{0\}$ 使 $\Gamma(Y^*,\pi^+(V))_{(s)}\neq 0$。若 $s>0$，取 $0\neq v\in\Gamma(Y^*,\pi^+(V))_{(s)}$，由前面的引理 (ii) 对 $0\le i\le n$ 有 $E\partial_iv=(s-1)\partial_iv$。显然所有 $\partial_iv=0$（$0\le i\le n$）不可能，因为那会蕴含 $Ev=0$。因此 $\Gamma(Y^*,\pi^+(V))_{(s-1)}\neq 0$，由 $s$ 上的递降归纳得矛盾。故 $s$ 必为负。此时若 $0\neq v\in\Gamma(Y^*,\pi^+(V))_{(s)}$，由前面的引理 (i) 对 $0\le i\le n$ 有 $Ey_iv=(s+1)y_iv$。若 $y_iv=0$，截面 $v$ 的支集含于 $Y^*$ 与第 $i$ 个坐标超平面的交中。由于所有坐标超平面与 $Y^*$ 的交为空，至少对某 $0\le i\le n$ 有 $y_iv\neq 0$。这意味着 $\Gamma(Y^*,\pi^+(V))_{(s+1)}\neq 0$，由 $s$ 上的归纳再次得矛盾。

由此可知 $\Gamma(Y,j_+(\pi^+(V)))=\Gamma(Y^*,\pi^+(V))=0$，且由于 $Y$ 是仿射簇，$j_+(\pi^+(V))=0$。这进而蕴含 $\pi^+(V)=0$，$V=0$。
\end{proof}

对任意拟凝聚 $\mathcal{D}_X$-模 $V$，存在自然态射 $\mu_V:\Delta(\Gamma(X,V))\to V$。显然 $\mu$ 是函子 $\Delta\circ\Gamma$ 到 $\mathcal{M}_{qc}(\mathcal{D}_X)$ 上恒等函子的自然变换。

\begin{lemma}
自然变换 $\mu$ 是函子 $\Delta\circ\Gamma$ 到 $\mathcal{M}_{qc}(\mathcal{D}_X)$ 上恒等函子的同构。
\end{lemma}

\begin{proof}
考虑拟凝聚 $\mathcal{D}_X$-模的正合列
$$0\longrightarrow K\longrightarrow\Delta(\Gamma(X,V))\longrightarrow V\longrightarrow C\longrightarrow 0.$$
由于 $\Gamma(X,-)$ 正合，得到 $D_X$-模的正合列
$$0\to\Gamma(X,K)\to\Gamma(X,\Delta(\Gamma(X,V)))\to\Gamma(X,V)\to\Gamma(X,C)\to 0.$$
由前面的引理中间箭头是同构，故 $\Gamma(X,K)=\Gamma(X,C)=0$。由此最终断定 $K=C=0$。
\end{proof}

这立即蕴含如下结果。

\begin{theorem}
函子 $\Gamma(X,-)$ 是范畴 $\mathcal{M}_{qc}(\mathcal{D}_X)$ 与 $\mathcal{M}(D_X)$ 的等价。其逆为 $\Delta$。
\end{theorem}

\begin{corollary}
$X=P^n$ 上任意拟凝聚 $\mathcal{D}_X$-模由其整体截面生成。
\end{corollary}

现在把这些结果推广到光滑仿射簇与射影空间的乘积。设 $X=P^n$，$Y$ 是光滑仿射簇，$\pi:X\times Y\to X$ 是自然投影。则 $\pi$ 是仿射态射。事实上若 $U\subset X$ 是开仿射子簇，$\pi^{-1}(U)=U\times Y$ 是 $X\times Y$ 的开仿射子簇。因此对任意拟凝聚 $\mathcal{O}_X$-模 $V$，
$$H^p(X\times Y,V)=H^p(X,\pi_*(V)),\quad p\in\mathbb{Z}.$$

若 $V$ 是 $\mathcal{D}_{X\times Y}$-模，$U\subset X$ 是开仿射子集，则
$$\pi_*(V)(U)=V(\pi^{-1}(U))=V(U\times Y)=\Gamma(U\times Y,V)$$
是 $D(U\times Y)$-模。由于 $D(U\times Y)=D(U)\otimes_kD(Y)$，$\pi_*(V)(U)$ 有由映射 $T\longmapsto T\otimes1$ 诱导的自然 $D(U)$-模结构。因此 $\pi_*(V)$ 有自然的 $\mathcal{D}_X$-模结构，此结构与 $\mathcal{O}_X$-模结构相容。由于 $\pi_*$ 保持拟凝聚性，若 $V$ 是拟凝聚 $\mathcal{D}_{X\times Y}$-模，直像 $\pi_*(V)$ 是拟凝聚 $\mathcal{D}_X$-模。故
$$H^p(X\times Y,V)=H^p(X,\pi_*(V))=0,\quad p>0,$$
且函子 $\Gamma(X\times Y,-)$ 在 $\mathcal{M}_{qc}(\mathcal{D}_{X\times Y})$ 上正合。另一方面若 $0=\Gamma(X\times Y,V)=\Gamma(X,\pi_*(V))$，由前面的结果 $\pi_*(V)=0$。若 $U\subset X$ 是开仿射子集，$0=\pi_*(V)(U)=V(U\times Y)$。由于 $U\times Y$ 是仿射簇且 $V$ 拟凝聚，可知 $V|_{U\times Y}=0$。由于 $U$ 任意，这意味着 $V=0$。

这证明前一定理的如下推广。

\begin{theorem}
设 $X=P^n$，$Y$ 是光滑仿射簇。则函子 $\Gamma(X\times Y,-)$ 是范畴 $\mathcal{M}_{qc}(\mathcal{D}_{X\times Y})$ 与 $\mathcal{M}(D_{X\times Y})$ 的等价。其逆为 $\Delta$。
\end{theorem}

特别地，拟凝聚 $\mathcal{D}_{X\times Y}$-模由其整体截面生成。